\documentclass[a4paper]{article}
\usepackage[pdftex]{graphicx}
\usepackage[utf8]{inputenc}
\usepackage{enumerate}
\usepackage{siunitx}
\sisetup{locale=DE} 
\usepackage{href-ul}
\hypersetup{
	colorlinks=true,
	linkcolor=blue,
	urlcolor=blue}
\usepackage{icomma}
\usepackage{amssymb}
\usepackage{geometry}
\geometry{a4paper, top=15mm, left=15mm, right=15mm, bottom=15mm,
	headsep=10mm, footskip=12mm}
%Source: img_0106

% Start the document
\begin{document}
%Titel
{\bf Schulaufgabe aus der Physik }\\
\begin{enumerate}[1.]
%Aufgabe 1
{\bf \item In einem Versuch wurde die Längenänderung eines Stabes in Abhängigkeit von der Temperatur gemessen.} Es ergab sich folgende Messreihe:

\begin{tabular}{cccccccc}
$\vartheta$ in $\SI{}{\celsius}$ & 0 & 10 & 15 & 20 & 28 & 35 & 40\\
$l$ in $\SI{}{\milli\meter}$ & 1000,0 & 1000,3 & 1000,6 & 1000,7 & 1000,9 & 1001,3 & 1001,5 
\end{tabular}

\begin{enumerate}[a)]
\item Werte den Versuch nun graphisch aus, indem du ein $\Delta l$ -- $\Delta \vartheta$ -- Diagramm zeichnest.
\item Werte die Messreihe numerisch aus und ermittle rechnerisch die Längen\-ausdehnungszahl 
$\alpha$.
\end{enumerate}

%Aufgabe 2
{\bf \item Ein Stahltank mit dem Volumen $V=60~l$ ist bei einer Temperatur von $\SI{15}{\celsius}$ vollständig mit Benzin gefüllt.} Wie viel Benzin läuft über, wenn sich die Temperatur auf $ \SI{40}{\celsius}$ erhöht?\\
$[ \gamma_{Benzin}= \SI{0,00106}{\kelvin^{-1}} \quad 
\gamma_{Stahl}= \SI{0,000036}{\kelvin^{-1}}$ 

%Aufgabe 3
{\bf \item Ein kaltes Flüssigkeitsthermometer wird schnell in heißes Wasser getaucht.}
 Man beobachtet, dass die Flüssigkeitssäule zuerst auf der Skala nach unten geht, um dann kehrtzumachen und anzusteigen, bis das Thermometer schließlich den tatsächlichen Temperaturwert anzeigt. Wie ist dieses anfängliche nach unten weichen zu erklären ?

%Aufgabe 4
\begin{minipage}{0.6\textwidth}
{\bf \item In einem geschlossenen Rohrsystems des in der Skizze ersichtlichen Querschnittes befindet sich Wasser der Temperatur $\SI{7}{\celsius}$.} Nun  wird der Schenkel A mit Eis umwickelt, sodass sich das Wasser dort zunächst auf $\SI{4}{\celsius}$ abkühlt. 
\begin{enumerate}[a)]
\item Erkläre nun, was zu beobachten ist.   
\item Was passiert, wenn sich das Wasser weiter bis  $\SI{1}{\celsius}$ abkühlt? 
\end{enumerate}
\end{minipage}
\hfill
\begin{minipage}{0.35\textwidth}
\includegraphics[width=4 cm]{zirk106.pdf}
\end{minipage}

%Aufgabe 5
{\bf \item Aufgabe}
\begin{enumerate}[a)]
\item wie lautet das Gesetz von Gay-Lussac über das Verhalten von  Volumen und Temperatur idealer Gase und unter welcher Voraussetzung ist es nur gültig?
(Das Gesetz ist in Worten und als Formel anzugeben)
\item Was besagt das Gesetz von Boyle-Mariotte über den  Druck  und das Volumen eines idealen Gases? Leite aus diesen beiden Gesetzen die allgemeine Gasgleichung her.
\end{enumerate}

{\bf \item Ein Wetterballon steigt vom Boden in eine Höhe von 10000 m auf.} Dabei gelten folgende Bedingungen: 

\begin{tabular}{ccc}
& Boden & Höhe \\
Temperatur & $\SI{18}{\celsius} $ & $\SI{-50}{\celsius}$ \\
Druck & $\SI{1000}{\hecto\pascal} $ & $\SI{264}{\hecto\pascal}$ \\
Volumen & $\SI{1,2}{\meter^3}$ & \\
\end{tabular}

Berechne das Volumen des Ballons in der Höhe.

{\bf \item Einem abgeschlossenen Gas wird mittels Strahlung Energie zugeführt. }
\begin{enumerate}[a)]
\item Zunächst bleibt das Volumen des Gases konstant. Was lässt sich über die Änderung der kinetischen Energie $E_{kin}$ und der potentiellen Energie $E_{pot}$ der Teilchen des Gases sagen?
\item Im weiteren Verlauf dehnt sich das Gas unter Beibehaltung seiner Temperatur aus (z. B. ein Kolben wird in einem Zylinder vom Gas nach außen geschoben). Was gilt jetzt für die Änderung der  $E_{kin}$ und der $E_{pot}$ der Teilchen?
\end{enumerate}

\end{enumerate}

\fbox{
	\begin{minipage}{0.5\textwidth}
		Zur Lösung bitte \href{https://www.okuyakl.de/phys/p9thdL106/ll106.pdf}{hier klicken} oder den QR-Code scannen.\\
	Weitere Arbeitsblätter gibt es unter 
	
	\href{https://www.okuyakl.de}{www.okuyakl.de}
	\end{minipage}
	\hfill
	\begin{minipage}{0.4\textwidth}
		\includegraphics[width=1.5 cm]{../../viecher/zwe03}
		\includegraphics[width=3 cm]{qr106}
		\includegraphics[width=2 cm]{../../viecher/afanticon1}
		
	\end{minipage}}

\end{document}